\documentclass[11pt]{article}
\usepackage[T1]{fontenc}
\usepackage[utf8]{inputenc}
\usepackage{lmodern,microtype}
\usepackage[margin=1in]{geometry}
\usepackage{amsmath,amssymb,amsthm,mathtools}
\usepackage{booktabs,array,longtable}
\usepackage{xcolor,listings,fancyhdr}
\usepackage[colorlinks=true,linkcolor=blue!45!black,urlcolor=blue!45!black,citecolor=blue!45!black]{hyperref}
\hypersetup{pdftitle={A smaller induction selector for iPRA},
  pdfauthor={Research note prepared for Vladimir Reshetnikov},
  pdfsubject={Unrestricted mirror-selected six-rule calculus and bi-interpretation with iPRA}}
\usepackage{enumitem}
\setlist{itemsep=3pt,topsep=5pt}
\setlength{\emergencystretch}{2em}
\setlength{\headheight}{14pt}
\pagestyle{fancy}\fancyhf{}
\fancyhead[L]{\small A smaller induction selector for iPRA}
\fancyhead[R]{\small Mirror core}
\fancyfoot[C]{\thepage}
\newtheorem{theorem}{Theorem}[section]
\newtheorem{lemma}[theorem]{Lemma}
\newtheorem{proposition}[theorem]{Proposition}
\newtheorem{corollary}[theorem]{Corollary}
\theoremstyle{definition}\newtheorem{definition}[theorem]{Definition}
\newcommand{\Mcal}{\mathcal M}
\newcommand{\Ncal}{\mathcal N}
\newcommand{\Scal}{\mathcal S}
\newcommand{\FV}{\operatorname{FV}}
\newcommand{\Mir}{\operatorname{Mir}}
\newcommand{\ipra}{\mathsf{iPRA}}
\newcommand{\ceq}{\mathrel{\equiv_{\rm syn}}}
\lstset{basicstyle=\ttfamily\footnotesize,columns=fullflexible,
  keepspaces=true,breaklines=true,showstringspaces=false,
  frame=single,rulecolor=\color{black!15},xleftmargin=4pt,xrightmargin=4pt,
  keywordstyle=\color{blue!45!black},commentstyle=\color{black!55}}
\title{A smaller induction selector for iPRA\\
  \large An unrestricted six-rule presentation with one mirrored iterator}
\author{Research note prepared for Vladimir Reshetnikov}
\date{November 11, 2026}
\begin{document}
\maketitle
\begin{abstract}
The current ConciseIPRA core already has six rule families on raw iterator
terms and unrestricted equality/$Q$/$D$ formulas. This note makes a modest
further simplification: replace its two-iterator diagonal test for
progressive introduction by the one-iterator pattern
$a(x,y)=I(y,a(x,y);u.h)$. The language, ordinary readings, and remaining
rules are unchanged. We prove equality of the complete contextual
consequence relations, and hence transfer both interpretation composites
and equiconsistency with full first-order intuitionistic primitive
recursive arithmetic. Arbitrary numerical substitution is admissible by
direct induction, using an inherited repair for newly selected bases.
Matched standalone Python checkers differ only in the selector: the
normalized source decreases from 3,125 to 3,109 bytes, from 1,022 to 1,012
significant tokens, and from 1,091 to 1,082 AST nodes. The gain is a
specialization and combination of existing ideas, not a new rule-count
reduction, global checker-size record, or absolute minimality theorem.
\end{abstract}
\tableofcontents
\newpage

\section{Result, baseline, and scope}

The comparison is pinned to ProveIt commit
\texttt{e1d9e3b12edd642ea2f5b572bb43febceaa50d93}, in
\texttt{Logic/PrimitiveRecursiveArithmetic/ConciseIPRA}. The current
README indexes 67 research packages and a Lean development. In particular,
\texttt{NativeSix.lean} defines the six-family calculus \texttt{SPrf},
proves \texttt{sprf\_iff\_nprf}, and proves substitution admissibility
\cite{readme,six}. That is the mathematical baseline used here.

\begin{theorem}[Main result]\label{thm:main}
Let $\Mcal$ be the calculus in Section~\ref{sec:spec}, and let $\Ncal$
be the current eight-constructor native kernel \texttt{NPrf}.
For every finite context $\Gamma$ and every native formula $A$,
\[
 \Gamma\vdash_{\Mcal}A
 \quad\Longleftrightarrow\quad
 \Gamma\vdash_{\Ncal}A.
\]
Both transformations are effective on finite derivations and leave the
sequent syntax unchanged. Consequently $\Mcal$ presents a theory
bi-interpretable and equiconsistent with full first-order $\ipra$.
Both numerical interpretation domains are whole; equality is actual
equality; and the two numerical composite isomorphisms have graph $x=y$.
\end{theorem}

The precise simplification is in the introduction selector:
\[
\begin{array}{ll}
\text{Current diagonal:}&I(y,y;u.h)=I(y,y;v.k),\\[2pt]
\text{Proposed mirror:}&a(x,y)=I(y,a(x,y);u.h).
\end{array}
\]
The latter recognizes one iterator skeleton and a repeated term. It also
allows a one-iterator lift of an arbitrary equation when recovering
induction. It does not restrict the grammar of public or intermediate
formulas. Neither a normalizer nor a primitive term-template parameter
is added to the checker.

\paragraph{Attribution.}
The mirror shape is already present in \texttt{Research/mirror\_guard},
where it is part of a restricted formula grammar and a different
transport interface \cite{mirror}. Two-case Eval is already present in
the current Lean development \cite{six}. The observation that a guard
need only be preserved under substitution, with newly true cases repaired
by a provable base, is explicitly established in
\texttt{Research/reflexive\_base/proof.tex}, including a raw-syntax
variant \cite{reflexive}. Those ingredients are inherited. The present
contribution is their particular combination on the full raw grammar,
the exact comparison with \texttt{SPrf}, and the supplied executable
evidence. No priority beyond the inspected snapshot is claimed.

\section{Complete specification of the proposed calculus}\label{sec:spec}

\subsection{Numerical terms and formulas}
There is one numerical sort. The raw term grammar is
\[
 t ::= x\mid 0\mid S(t)\mid I(n,a;u.h).
\]
Only the state variable $u$ is bound in $h$. The counter $n$ and
initializer $a$ are outside that binder. An update can have fixed
numerical parameters. Terms are finite syntax trees, not equivalence
classes. Substitution is ordinary capture-free substitution, with no
normalization step.

The formula grammar is unrestricted:
\[
 A ::= (s=t)\mid Q_y A\mid D_{x,y}A.
\]
The variables bound by $D$ are distinct. Write $B(t,r)$ for simultaneous
capture-free replacement of its distinguished free coordinates $x,y$ by
$t,r$. The designated contradiction is $0=S0$. There are no primitive
ordinary connectives, falsity symbol, predicate symbols other than
equality, or additional numerical operations.

For orientation, the ordinary readings are
\[
 Q_yB(y) := B(0)\to\forall y\,B(y),
 \qquad
 D_{x,y}B(x,y) :=
 \exists x\bigl(B(x,0)\to\forall y\,B(x,y)\bigr).
\]
These are interpretation clauses, not extra axioms of the native
calculus. Contexts are finite lists of reusable formulas. Their order
and multiplicities have no logical effect, although a checker may use
indices to identify a hypothesis.

\subsection{The selector}
Define $\Mir(B)$ to hold exactly when the \emph{uninstantiated} matrix
$B(x,y)$ is literally
\[
             a(x,y)=I(y,a(x,y);u.h(x,y,u)).             \tag{M}
\]
The counter must be the distinguished bound variable $y$. The
initializer must be the same raw term tree as the left side of the
equation. The body $h$ is arbitrary, and its state binder is fresh.
Define
\[
 \eta_B(z)=\begin{cases}z,&\Mir(B),\\0,&\neg\Mir(B).\end{cases}
\]
This is a syntactic test. It asks for neither arithmetic equality nor
provability. It is performed before substituting the existential
witness. For example $x=I(y,0;u.0)$ is not selected, even when the
subsequently supplied witness is $0$.

\subsection{All primitive rules}
The following six families are the entire calculus:
\begin{align*}
 &\frac{A\in\Gamma}{\Gamma\vdash A}
       &&\mathrm{Hyp},\\[5pt]
 &\frac{}{\Gamma\vdash t=t}
       &&\mathrm{Refl},\\[5pt]
 &\frac{\Gamma\vdash a=t\qquad
         \Gamma\vdash I(0,a;u.h)=r}
        {\Gamma\vdash r=t}
       &&\mathrm{Eval}_0,\\[5pt]
 &\frac{\Gamma\vdash I(c,a;u.h)=t\qquad
         \Gamma\vdash I(Sc,a;u.h)=r}
        {\Gamma\vdash r=h[t/u]}
       &&\mathrm{Eval}_S,\\[5pt]
 &\frac{\Gamma\vdash Q_yB(y)\qquad\Gamma\vdash B(0)}
        {\Gamma\vdash B(t)}
       &&\mathrm{QE},\\[5pt]
 &\frac{\Gamma\vdash D_{x,y}B(x,y)\qquad
         \Gamma,Q_yB(b,y)\vdash A}
        {\Gamma\vdash A}
       &&\mathrm{DE},\\[5pt]
 &\frac{\Gamma,B(t,\eta_B(z))\vdash B(t,Sz)}
        {\Gamma\vdash D_{x,y}B(x,y)}
       &&\mathrm{DI}_{M}.
\end{align*}
The two Eval rows are the two cases of one family, exactly as in the
baseline's six-rule count. There are seven inference constructors if
those cases are counted separately. This convention is explicit in both
systems; no rule has been hidden by a change of counting convention.

In DE, $b$ is fresh for $\Gamma$, $A$, and every passive parameter of
$B$. In DI, $z$ is fresh for $\Gamma$, $t$, and every passive parameter
of $B$. The displayed local assumptions are discharged. All substitutions
avoid capture. There are no other logical or nonlogical axioms, no
primitive structural rules, no successor axioms, no separate computation
axioms, and no general induction rule. All of these needed operations
are derived below.

\section{Structural operations, equality, and computation}

\subsection{The order of structural arguments}
\begin{lemma}\label{lem:struct}
Weakening, injective renaming of numerical variables, and simultaneous
ordinary-assumption cut are admissible in $\Mcal$, without using
general numerical substitution or equation induction.
\end{lemma}
\begin{proof}
Weakening is induction on derivations. An injective renaming preserves
constructor shapes and is injective on raw term trees. It therefore
preserves and reflects both the distinguished counter test and equality
of the two mirrored occurrences of $a$. It preserves every primitive
rule, including the DI mode. Renaming a derivation into a larger fresh
scope is a special case.

For cut, induct on the derivation into which the supplied proofs are
inserted. Replace a hypothesis leaf by its supplied derivation. At a
binder rule, rename the eigenvariable away from the finitely many
inserted data, lift the supplied derivations by the injective renaming
just proved, and retain the local assumption by Hyp. The matrix of DI
is unaffected by replacement of hypotheses, so its mode is unchanged.
This proves simultaneous cut and its one-assumption specialization.
\end{proof}

Exchange and contraction are immediate consequences of reusable
hypotheses and weakening. General numerical substitution is deliberately
postponed: a noninjective substitution can make a previously false
mirror test true.

\subsection{Explicit equality derivations}
All the following arguments use Refl and Eval alone, so they are the
same elementary bootstrap as in \texttt{NativeSix} \cite{six}.

\begin{lemma}[Joining equations]\label{lem:join}
From $s=t$ and $s=r$ derive $r=t$.
\end{lemma}
\begin{proof}
Put $J=I(0,s;u.u)$. Eval$_0$ on $s=r$ and $J=J$ gives $J=r$.
Eval$_0$ on $s=t$ and this $J=r$ gives $r=t$.
\end{proof}

Symmetry follows by joining $s=s$ and $s=t$. To obtain transitivity
from $s=t$ and $t=u$, reverse the first equation and join the two
equations with common left side $t$.

Both computation equations are derived:
\[
 I(0,a;u.h)=a,
 \qquad I(Sc,a;u.h)=h[I(c,a;u.h)/u].             \tag{I0, IS}
\]
Use the appropriate Eval rule on reflexivities of its two sources.

\begin{lemma}[One-term transport]\label{lem:transport}
For any fixed term template $h(v)$, from $s=t$ and $h(s)=r$ derive
$r=h(t)$. Consequently all numerical congruence and equational
replacement rules are available.
\end{lemma}
\begin{proof}
Let $J=I(0,s;v.h(v))$ and $K=I(S0,s;v.h(v))$.
I0 gives $J=s$, hence $J=t$ by transitivity. Eval$_S$ on $J=s$ and
$K=K$ gives $K=h(s)$, hence $K=r$. Eval$_S$ on $J=t$ and $K=r$
now gives $r=h(t)$. Choosing the second premise as a reflexivity and
using symmetry yields ordinary congruence. Transport in both sides of
an equation is two applications with the corresponding templates.
\end{proof}

\section{Recovering cases before induction}

\begin{lemma}[Provable mirror bases]\label{lem:base}
If $\Mir(B)$, then $\Gamma\vdash B(t,0)$ for every context and witness.
\end{lemma}
\begin{proof}
The instance is $a(t,0)=I(0,a(t,0);u.h(t,0,u))$. This is the symmetry
of I0. No induction, coverage, or logical macro is used.
\end{proof}

\begin{lemma}[All case introductions]\label{lem:allcases}
For every raw matrix $B$, the following rule is admissible:
\[
 \frac{\Gamma,B(t,0)\vdash B(t,Sz)}
      {\Gamma\vdash D_{x,y}B}.
\]
\end{lemma}
\begin{proof}
If the matrix is not mirrored, this is primitive DI. If it is mirrored,
cut the proof from Lemma~\ref{lem:base} into the supplied branch, obtaining
$B(t,Sz)$ under $\Gamma$ alone. Weaken by $B(t,z)$ and use progressive
DI. The step variable remains fresh because the removed base contains
no $z$.
\end{proof}

\begin{lemma}[Case-closure introduction and numerical cases]\label{lem:ncases}
For every native formula $F$, both rules are admissible:
\[
 \frac{\Gamma,F(0)\vdash F(Sz)}{\Gamma\vdash Q_yF(y)},
 \qquad
 \frac{\Gamma\vdash F(0)\qquad\Gamma\vdash F(Sz)}
      {\Gamma\vdash F(v)}.
\]
In the second rule $z$ is fresh for the context, passive parameters,
and requested target $v$.
\end{lemma}
\begin{proof}
Introduce $D_{w,y}F(y)$ by the derived case rule with an unused witness
coordinate $w$ and witness $0$. DE exposes $Q_yF(y)$, independent of
the fresh witness, so Hyp returns that closure and DE closes. This is
the first rule. For the second, weaken the successor premise by $F(0)$,
apply the first rule, and apply QE to the supplied base at $v$.
\end{proof}

The case rule is available for arbitrary compound formulas. Its proof
uses no equation induction. In particular, the numerical-case arguments
in the next section do not assume the induction they are used to recover.

\section{Recovering unrestricted equation induction}\label{sec:induction}

Write $C(c,a,b)=I(c,a;u.b)$, with $u$ unused. Computation gives
$C(0,a,b)=a$ and $C(Sv,a,b)=b$. Let
\[
 E(y)\equiv s(y)=r(y),\qquad
 G(y)\equiv s(y)=C(y,s(y),r(y)).                      \tag{Lift}
\]
The latter is a mirrored equation with one iterator. It has a provable
zero instance and its successor instance is interderivable with that
of $E$, using only computation and equality.

\begin{lemma}[Equational transfer]\label{lem:transfer}
For every target $v$, from $E(0)$ and $G(v)$ derive $E(v)$ without
equation induction.
\end{lemma}
\begin{proof}
Numerical cases establish the following unconditional equations:
\begin{align}
 C(v,s(0),s(v))&=s(v),                               \label{eq:Ls}\\
 C(v,r(0),C(v,s(v),r(v)))&=r(v).                     \label{eq:Lr}
\end{align}
At zero both left sides compute to the displayed zero seed. At a
successor the first equation uses one constant-update computation and
the second uses two. Apply Lemma~\ref{lem:ncases} to these equations
\emph{before} retaining any assumption depending on $v$. Then weaken
the proved equations into the desired context. This order is necessary
for the eigenvariable condition.

Under $G(v)$, congruence for the template $C(v,s(0),-)$ gives
\[
 C(v,s(0),s(v))=C(v,s(0),C(v,s(v),r(v))).
\]
The retained base $s(0)=r(0)$ changes the outer seed on the right.
Using \eqref{eq:Ls} and \eqref{eq:Lr}, the complete chain is
\begin{align*}
 s(v)
 &=C(v,s(0),s(v))\\
 &=C(v,s(0),C(v,s(v),r(v)))\\
 &=C(v,r(0),C(v,s(v),r(v)))\\
 &=r(v).
\end{align*}
Each step is a previously available equational inference.
\end{proof}

\begin{theorem}[Contextual equation induction]\label{thm:induction}
For every literal equation $E$ and arbitrary retained context $\Gamma$,
\[
 \frac{\Gamma\vdash E(0)\qquad\Gamma,E(z)\vdash E(Sz)}
      {\Gamma\vdash E(t)}
\]
is admissible, with the usual fresh $z$.
\end{theorem}
\begin{proof}
Use the lift $G$ above. Under $\Gamma,G(z)$, the base proof and
Lemma~\ref{lem:transfer} give $E(z)$. Cut this into the supplied step
derivation to get $E(Sz)$. The constant-update computation converts this
to $G(Sz)$. Thus the progressive premise for $G$ has been obtained.

With an unused existential coordinate $w$ and witness $0$, primitive
DI gives $D_{w,y}G(y)$. DE exposes $Q_yG(y)$. The theorem $G(0)$
and QE at $t$ yield $G(t)$; the retained base and
Lemma~\ref{lem:transfer} yield $E(t)$. Close DE. No substitution into
the supplied step derivation was needed, apart from injective fresh
renaming of its eigenvariable.
\end{proof}

\begin{corollary}[Native equation-progress rule]\label{cor:qprogress}
From $\Gamma,E(z)\vdash E(Sz)$ derive $\Gamma\vdash Q_yE(y)$.
\end{corollary}
\begin{proof}
Use the first rule of Lemma~\ref{lem:ncases}. In its branch retain
$E(0)$ and apply Theorem~\ref{thm:induction} at target $Sv$ to obtain
the required $E(Sv)$. The step eigenvariable is renamed fresh from
this outer $v$ and all retained data. Then discharge $E(0)$ by the
case-closure rule.
\end{proof}

\section{Exact sequent equivalence and numerical substitution}

\subsection{Proof of the main equivalence}
For $\Mcal\to\Ncal$, Hyp, Refl, Eval, QE, and DE have the existing
native derivations. A nonmirrored DI is the native case introduction.
A mirrored DI has a literal equation matrix. The native equation-progress
rule gives $Q_yB(t,y)$; under $B(t,0)$, QE supplies $B(t,Sz)$, and
native case introduction with witness $t$ gives $D_{x,y}B$. This
translates every primitive rule without changing the concluding sequent.

For $\Ncal\to\Mcal$, derive I0, IS, and one-term transport as in
Section~3. Hyp, QE, and DE are primitive. Replace native case introduction
by Lemma~\ref{lem:allcases} and native equation-progress introduction by
Corollary~\ref{cor:qprogress}. Induction on the native derivation completes
the transformation. Together these prove the first assertion of
Theorem~\ref{thm:main}. The existing equality of \texttt{SPrf} and
\texttt{NPrf} gives exact equality with the six-rule baseline as well.

\subsection{Substitution may change the mode}
\begin{theorem}[Direct substitution admissibility]\label{thm:subst}
If $\Gamma\vdash_{\Mcal}A$ and $\sigma$ is any capture-free substitution
of terms for passive numerical variables, then
$\Gamma\sigma\vdash_{\Mcal}A\sigma$.
\end{theorem}
\begin{proof}
The mirror test is preserved by substitution: the repeated initializer
and left side receive the same substitution, while the distinguished
bound counter remains $y$. It need not be reflected. For example,
$p=I(y,q;u.0)$ becomes mirrored after both $p,q$ are replaced by $0$.

Induct on the derivation. Rename each eigenvariable fresh from the finite
union of free variables in the relevant substitution images. All rules
other than DI commute with the corresponding lifted substitutions. If
DI originally used progress, preservation gives the same mode after
substitution, so apply the induction hypothesis and primitive DI. If
DI originally used cases, the induction hypothesis gives a substituted
case premise; apply the already derived unrestricted case introduction
of Lemma~\ref{lem:allcases}. That rule repairs a newly mirrored case by
cutting its provable zero base and weakening to progress.
\end{proof}

This proof is internal to the new calculus; it does not translate to
$\Ncal$ and back. The repair principle is inherited from
\texttt{reflexive\_base} \cite{reflexive}. Unlike the current diagonal
selector, the mirror selector is not preserved \emph{and reflected} by
arbitrary substitutions. The theorem claims admissibility with an
explicit proof repair, not that the primitive DI instance always keeps
the same mode.

% ORDINARY_BRIDGE_INSERT

\section{Checker, exact comparison, and verification}\label{sec:implementation}

\subsection{Public certificate contract}
The standalone checker uses Python 3.10 or later and only the standard
\texttt{json} and \texttt{sys} modules. The measured runtime is CPython
3.12.14. Input is one JSON object with exactly four fields:
\begin{lstlisting}
{"variables": n, "context": [A0, A1, ...], "goal": A, "proof": p}
\end{lstlisting}
Here $n$ is an exact nonnegative integer, not a Boolean. Terms use
\texttt{["v",i]}, \texttt{["z"]}, \texttt{["s",t]}, and
\texttt{["i",n,a,h]}; formulas use \texttt{["=",s,t]},
\texttt{["q",B]}, and \texttt{["d",B]}. Binding is de Bruijn:
the distinguished $y$ is index $0$ and $x$ index $1$ in a $D$ matrix.
The iterator increments binding depth only in its update body.

The proof tags are \texttt{hyp}, \texttt{refl}, \texttt{eval},
\texttt{qe}, \texttt{de}, and \texttt{di}. Tag fields and every supplied
term and formula are checked at the appropriate scope; malformed active
syntax, Boolean variable indices, invalid references, and escaping
eigenvariables are rejected. Each recursive inference computes its
conclusion and the top-level goal is checked against that result.
The only successful public output is \texttt{ACCEPT} followed by a
newline and exit status zero. A rejected input exits unsuccessfully.

The checker is a direct implementation of Section~\ref{sec:spec}.
Validation of terms/formulas and a structural induction on successful
inference calls establish that accepted certificates are derivations of
the specified calculus. Conversely, every finite derivation has the
corresponding tree certificate. The completeness statement, like that
of the repository checkers, assumes unbounded available resources;
a fixed Python recursion or memory limit may reject a very large valid
certificate. Such a failure is never acceptance.

\subsection{Matched measurements}
The comparator implements the current \texttt{SPrf} calculus in the
same typed Python framework. It is a supplied matched implementation,
not a claim that this exact Python file already appears upstream.
Both files have the same parser, term and formula validators,
capture-avoiding substitution, context interface, and all other rule
branches. The only executable-source difference is the DI guard.

\begin{center}
\begin{tabular}{lrrr}
\toprule
Complete normalized module & Bytes & Tokens & AST nodes\\
\midrule
Current diagonal selector & 3125 & 1022 & 1091\\
Proposed mirror selector  & 3109 & 1012 & 1082\\
Decrease                 &   16 &   10 &    9\\
\bottomrule
\end{tabular}
\end{center}

The measurement removes genuine docstrings, applies \texttt{ast.unparse},
and appends one line feed. Bytes mean normalized UTF-8 bytes. Significant
tokens exclude the end marker, encoding, comments, newlines, indentation,
and dedentation. AST nodes are counted after docstring removal. The
archive contains the complete source, measurement code, normalized
copies, and a one-line executable diff.

This is a 0.51\% byte reduction against the matched six-rule control.
It is not a reduction below the repository's 516-byte compact codec,
nor a new minimum among all of its public-goal or fully validated
interfaces. Those smaller programs use different interfaces and
encodings. A six-rule presentation and a shorter matched recognizer
are the supported simplicity claims.

The tradeoff is also concrete: the diagonal guard inspects fixed
constructor/variable positions, whereas the mirror guard may compare
two large term trees. That comparison is linear in their size in the
worst case. The source reduction does not imply uniformly faster
checking or shorter proofs for every theorem.

\subsection{Evidence and limits}
% VALIDATION_INSERT

The general equivalence and interpretation claims are mathematical
proofs in this note. Executed certificates support the implementation
and selected proof transformations; finite testing cannot establish
the universal metatheorems. The existing Lean source endpoints were
read at the pinned revision. Their project was not rebuilt here, and
there is no new Lean or Rocq formalization of the mirror calculus or
of the Python checker. The included Lean files are reference excerpts
with their original imports, not a standalone formalization project.

\section{Further research questions}
\begin{enumerate}
\item Can the same selector be integrated into the repository's smallest
complete certificate interface with a net source reduction after all
validation and scope checks are counted? The present comparison does
not answer that question.
\item What is the smallest family of syntactic base recognizers that
supports an equational decoder for every equation? The mirror and
diagonal shapes give two concrete instances, but neither establishes
a lower bound.
\item Can the mirror comparison be made constant-time by syntax sharing
without moving a larger validation obligation into the trusted input
reader? Any complexity comparison must count the representation and
its integrity checks.
\item Can the equation-induction translator be formalized directly in
Lean with fewer proof dependencies than the current diagonal
construction? The two unconditional selector equations are a useful
target for such a formalization.
\item Can one prove a substantive lower bound for a precisely specified
class of complete proof checkers? Absolute claims over unrestricted
programming languages or encodings are not supported by the present
experiments.
\end{enumerate}

\appendix
\section{Complete standalone checker}\label{app:checker}
The following is the complete selected acceptance module. No proof
generator, source translator, reference theorem, or measurement helper
is called by it.
% CHECKER_INSERT

\clearpage
\section*{Inherited arithmetic appendices}
\addcontentsline{toc}{section}{Inherited arithmetic appendices}
The following three sections reproduce the pinned ConciseIPRA
predecessor, equality-reflection, and primitive-recursion arguments
\cite{erasure,forward,progress}. They are inherited foundations rather
than new reductions. They make the dependency of the iPRA bridge
reviewable without relying on standard-model agreement. The original
source files and their license are also included in the archive.
% ARITHMETIC_INSERT

\clearpage
\begin{thebibliography}{9}
\raggedright
\bibitem{readme} V. Reshetnikov, ProveIt, ConciseIPRA README, pinned commit
\texttt{e1d9e3b12edd642ea2f5b572bb43febceaa50d93}.
\href{https://github.com/VladimirReshetnikov/ProveIt/blob/e1d9e3b12edd642ea2f5b572bb43febceaa50d93/Logic/PrimitiveRecursiveArithmetic/ConciseIPRA/README.md}{Repository README}.
\bibitem{six} ProveIt, \texttt{Lean/ConciseIPRA/NativeSix.lean} and
\texttt{NativeDiagonal.lean}, same pinned commit.
\href{https://github.com/VladimirReshetnikov/ProveIt/blob/e1d9e3b12edd642ea2f5b572bb43febceaa50d93/Logic/PrimitiveRecursiveArithmetic/ConciseIPRA/Lean/ConciseIPRA/NativeSix.lean}{Six-rule calculus and its theorems}.
\bibitem{mirror} ProveIt, \texttt{Research/mirror\_guard}, same commit.
\href{https://github.com/VladimirReshetnikov/ProveIt/tree/e1d9e3b12edd642ea2f5b572bb43febceaa50d93/Logic/PrimitiveRecursiveArithmetic/ConciseIPRA/Research/mirror_guard}{One-sided computation guard package}.
\bibitem{reflexive} ProveIt, \texttt{Research/reflexive\_base/proof.tex},
same commit.
\href{https://github.com/VladimirReshetnikov/ProveIt/blob/e1d9e3b12edd642ea2f5b572bb43febceaa50d93/Logic/PrimitiveRecursiveArithmetic/ConciseIPRA/Research/reflexive_base/proof.tex}{Reflexive-base introduction and substitution repair}.
\bibitem{bridge} ProveIt, \texttt{Lean/ConciseIPRA/NativeBiInterpretation.lean},
same commit.
\href{https://github.com/VladimirReshetnikov/ProveIt/blob/e1d9e3b12edd642ea2f5b572bb43febceaa50d93/Logic/PrimitiveRecursiveArithmetic/ConciseIPRA/Lean/ConciseIPRA/NativeBiInterpretation.lean}{Faithful readings, both round trips, and equiconsistency}.
\bibitem{erasure} ProveIt, \texttt{Research/certified\_zero/reference/}\allowbreak
\texttt{scalar\_predecessor.tex}, same commit. Reproduced in the arithmetic appendix.
\bibitem{forward} ProveIt, \texttt{Research/certified\_zero/reference/}\allowbreak
\texttt{equality\_reflection.tex}, same commit. Reproduced in the arithmetic appendix.
\bibitem{progress} ProveIt, \texttt{Research/certified\_zero/reference/}\allowbreak
\texttt{primitive\_recursion.tex}, same commit. Reproduced in the arithmetic appendix.
\end{thebibliography}
\end{document}
